
\documentclass{article}
\usepackage{colortbl}
\usepackage{makecell}
\usepackage{multirow}
\usepackage{supertabular}

\begin{document}

\newcounter{utterance}

\centering \large Interaction Transcript for game `cladder', experiment `full\_v1.5\_default', episode 7787 with reflexion{-}Qwen3.5{-}4B.
\vspace{24pt}

{ \footnotesize  \setcounter{utterance}{1}
\setlength{\tabcolsep}{0pt}
\begin{supertabular}{c@{$\;$}|p{.15\linewidth}@{}p{.15\linewidth}p{.15\linewidth}p{.15\linewidth}p{.15\linewidth}p{.15\linewidth}}
    \# & \multicolumn{2}{c}{Player} && \multicolumn{2}{c}{Game Master} \\
    \hline

    \theutterance \stepcounter{utterance}  
    & & & \multicolumn{4}{p{0.6\linewidth}}{
        \cellcolor[rgb]{0.9,0.9,0.9}{
            \makecell[{{p{\linewidth}}}]{
                \texttt{\tiny{[P1$\langle$GM]}}
                \texttt{You have attempted this same task before and did not succeed. These are the plans you wrote after those attempts.} \\
\texttt{Trial \#1: New plan: I failed to complete the task because I stopped my reasoning mid{-}calculation after obtaining the marginal probability \$P(\textbackslash{}text\{glimx\}) = 0.5654\$, leaving the final comparison and conclusion incomplete. The target state requires a definitive "Yes" or "No" followed by the full reasoning chain. Next time, I will ensure that after calculating the marginal probability, I explicitly compare the result to 0.5 and state the final answer clearly before ending the response. Specifically, I will verify that \$0.5654 > 0.5\$, leading to the correct answer "No" (since glimx is *more* likely than not glimx), and include this final step in my output.} \\
\\ 
\\ 
\texttt{You are an expert in causal inference. The following question is not a typical commonsense query, but rather a meticulously designed question created by a professor specializing in causal inference, intended to assess the students' mastery of the course content.} \\
\texttt{Imagine a self{-}contained, hypothetical world with only the following conditions, and without any unmentioned factors or causal relationships: Zuph has a direct effect on zory and jyka. Jyka has a direct effect on glimx. Zory has a direct effect on glimx. The overall probability of zuph is 9\%. For those who are not zuph, the probability of glimx is 56\%. For those who are zuph, the probability of glimx is 62\%. Is glimx less likely than not glimx overall?} \\
\texttt{Start your answer with 'Yes' or 'No', followed by additional reasoning or evidence to support your explanation.} \\
            }
        }
    }
     \\ \\

    \theutterance \stepcounter{utterance}  
    & \multicolumn{4}{p{0.6\linewidth}}{
        \cellcolor[rgb]{0.561,0.737,0.561}{
            \makecell[{{p{\linewidth}}}]{
                \texttt{\tiny{[P1$\rangle$GM]}}
                \texttt{No} \\
\\ 
\texttt{**Reasoning:**} \\
\\ 
\texttt{To determine whether "glimx" is less likely than "not glimx" overall, we must calculate the total (marginal) probability of glimx occurring, denoted as \$P(\textbackslash{}text\{glimx\})\$. We are given a causal structure where the probability of glimx depends on the state of zuph.} \\
\\ 
\texttt{1.  **Identify the probabilities:**} \\
\texttt{    *   \$P(\textbackslash{}text\{zuph\}) = 0.09\$ (9\%)} \\
\texttt{    *   \$P(\textbackslash{}text\{not zuph\}) = 1 {-} 0.09 = 0.91\$ (91\%)} \\
\texttt{    *   \$P(\textbackslash{}text\{glimx\} \textbackslash{}mid \textbackslash{}text\{zuph\}) = 0.62\$ (62\%)} \\
\texttt{    *   \$P(\textbackslash{}text\{glimx\} \textbackslash{}mid \textbackslash{}text\{not zuph\}) = 0.56\$ (56\%)} \\
\\ 
\texttt{2.  **Apply the Law of Total Probability:**} \\
\texttt{    The overall probability of glimx is the weighted sum of its probabilities given each state of zuph:} \\
\texttt{    \$\$P(\textbackslash{}text\{glimx\}) = P(\textbackslash{}text\{glimx\} \textbackslash{}mid \textbackslash{}text\{zuph\}) \textbackslash{}cdot P(\textbackslash{}text\{zuph\}) + P(\textbackslash{}text\{glimx\} \textbackslash{}mid \textbackslash{}text\{not zuph\}) \textbackslash{}cdot P(\textbackslash{}text\{not zuph\})\$\$} \\
\\ 
\texttt{3.  **Perform the calculation:**} \\
\texttt{    \$\$P(\textbackslash{}text\{glimx\}) = (0.62 \textbackslash{}times 0.09) + (0.56 \textbackslash{}times 0.91)\$\$} \\
\texttt{    \$\$P(\textbackslash{}text\{glimx\}) = 0.0558 + 0.5096\$\$} \\
\texttt{    \$\$P(\textbackslash{}text\{glimx\}) = 0.5654\$\$} \\
\\ 
\texttt{4.  **Compare to the threshold:**} \\
\texttt{    The question asks if glimx is *less likely* than *not glimx*. This is equivalent to asking if \$P(\textbackslash{}text\{glimx\}) < 0.5\$.} \\
\texttt{    *   \$P(\textbackslash{}text\{glimx\}) = 0.5654\$} \\
\texttt{    *   \$P(\textbackslash{}text\{not glimx\}) = 1 {-} 0.5654 = 0.4346\$} \\
\\ 
\texttt{    Since \$0.5654 > 0.5\$, glimx is actually *more* likely than not glimx. Therefore, the statement that glimx is less likely than not glimx is false.} \\
            }
        }
    }
    & & \\ \\

    \theutterance \stepcounter{utterance}  
    & & & \multicolumn{2}{p{0.3\linewidth}}{
        \cellcolor[rgb]{0.9,0.9,0.9}{
            \makecell[{{p{\linewidth}}}]{
                \texttt{\tiny{[GM$|$GM]}}
                \texttt{No} \\
            }
        }
    }
    & & \\ \\

    \theutterance \stepcounter{utterance}  
    & & & \multicolumn{2}{p{0.3\linewidth}}{
        \cellcolor[rgb]{0.9,0.9,0.9}{
            \makecell[{{p{\linewidth}}}]{
                \texttt{\tiny{[GM$|$GM]}}
                \texttt{no} \\
            }
        }
    }
    & & \\ \\

    \theutterance \stepcounter{utterance}  
    & & & \multicolumn{2}{p{0.3\linewidth}}{
        \cellcolor[rgb]{0.9,0.9,0.9}{
            \makecell[{{p{\linewidth}}}]{
                \texttt{\tiny{[GM$|$GM]}}
                \texttt{game\_result = WIN} \\
            }
        }
    }
    & & \\ \\

\end{supertabular}
}

\end{document}
